\chapter{User Interface Design}
\ChapterLabel{UserInterfaceDesign}
\section{Chapter Responsibility}
\label{Section:UserInterfaceDesignResponsibility}
Define how people interact with the system and understand its state.

\noindent\textbf{Expected result:} An interaction design that supports stakeholder tasks and acceptance.

\section{Revision 0.1 Scope Notice}
\label{Section:DeferredUserInterfaceDesign}
\begin{fsaNote}[Deferred project content --- template reference only]
This chapter is retained in the complete document as template guidance.
Its sample artifacts, IDs, priorities, statuses, versions, outcomes, and
approval examples do not describe the FTDI245R project. In particular, any
sample Done, Pass, or release entry is not implementation, verification,
validation, or approval evidence for this project. Detailed project content
is deferred under \hyperref[Section:ProjectScope]{the revision scope}.
See \hyperref[Table:ProjectGapReview]{the reviewed gap list} for actual
decisions and unresolved items.
\end{fsaNote}

\section{How to Use This Chapter}
Describe personas, primary screens or workflows, interaction rules, feedback, accessibility, and essential states. Include wireframes or prototypes where they resolve meaningful ambiguity.
\section{Example}
Persona: a busy operations coordinator needs to submit and track requests quickly. The request screen groups required fields, identifies inline validation errors, confirms successful submission, supports keyboard operation, and exposes empty, loading, error, and completed states.
